\subsection{Successful-interaction skeleton and patches}
\label{sec:patches}

We now choose a coarse~$\sigma$-algebra for the graphical construction on which
to condition. It should determine every interaction between distinct sites
while leaving local randomness unobserved. An interaction between sites occurs
exactly when a mark with nonempty target rings at an active source. To specify
such an interaction, it is enough to retain its source, time, and target. We do
not retain its kind~$\alpha$: for fixed~$(i,t,S)$, a birth and a split involve
the same source and target sites and differ only in the resulting state of the
source. We also omit deaths, whose target is empty, and rings at inactive
sources, which do not change the process at all. The retained records form the
successful-interaction skeleton. Once this skeleton is fixed, its interaction
times divide each site line into consecutive intervals, called patches. The
omitted deaths and inactive-source rings occur within individual site lines
between these interaction times, while the omitted kind of a retained
interaction affects only its source line.

\begin{definition}[Successful interaction and its record]
\label{def:successful-interaction}
The deterministic initial interaction~$\iota_0$ from
\eqref{eq:initial-interaction} is declared successful and has record
\begin{equation*}
(\infty,0,A_0).
\end{equation*}
An interaction~$(i,t,\alpha,S)\in I^{\mathrm P}$ is
\emph{successful} when
\begin{equation*}
S\neq\varnothing
\qquad\text{and}\qquad
i\in A_{t-}.
\end{equation*}
Its successful-interaction record is~$(i,t,S)$. Thus the record remembers the
source, time, and target of the interaction, but not its kind~$\alpha$.
\end{definition}

For~$T<\infty$, let~$\mathcal I_T$ be the successful-interaction skeleton up to
time~$T$:
\begin{equation*}
\mathcal I_T
=
\cb{(\infty,0,A_0)}
\cup
\cb{(i,t,S):t\leq T\text{ and, for some }\alpha\in\cb{\beta,\delta},
(i,t,\alpha,S)\text{ is successful}}.
\end{equation*}
Thus~$\mathcal I_T$ contains exactly the retained records described above. Set
\begin{equation*}
\mathcal I=\bigcup_{T<\infty}\mathcal I_T,
\qquad
\mathcal G_T=\sigma\bb{Y_0,\mathcal I_T}.
\end{equation*}
Thus~$\mathcal G_T$ records the initial state and the successful-interaction
skeleton up to time~$T$. The local-finiteness statement in
Section~\ref{subsec:dual-graphical-construction} implies that~$\mathcal I_T$ is
finite almost surely.

We regard a record~$(j,t,S)\in\mathcal I$ as incoming at every site in~$S$ and
as outgoing at its source. On each site line, these interactions divide time
into consecutive intervals. Each interval begins with an incoming or outgoing
interaction and ends at the next successful interaction involving that site, if
there is one.

For~$i\in\Lambda$, let
\begin{align*}
\mathcal T^{\mathsf I}(i)
&=
\cb{t\geq0:(j,t,S)\in\mathcal I
  \text{ for some }j,S\text{ with }i\in S},
\\
\mathcal T^{\mathsf O}(i)
&=
\cb{t>0:(i,t,S)\in\mathcal I\text{ for some }S},
\\
\mathcal T(i)
&=
\mathcal T^{\mathsf I}(i)\cup\mathcal T^{\mathsf O}(i).
\end{align*}
For~$s\in\mathcal T(i)$, the next successful interaction involving~$i$ occurs
at
\begin{equation*}
e_i(s)=\inf\cb{u>s:u\in\mathcal T(i)},
\qquad
\inf\varnothing=\infty.
\end{equation*}

\begin{definition}[Patch]
\label{def:patch}
Let~$(j,s,S)\in\mathcal I$ be a record for which~$i=j$ or~$i\in S$. The patch
beginning at~$(i,s)$ is
\begin{equation*}
P=
\bb{i(P),[s(P),e(P)),\mathsf X(P)\mathsf Y(P),S(P)},
\end{equation*}
where
\begin{equation*}
i(P)=i,
\qquad
s(P)=s,
\qquad
e(P)=e_i(s),
\qquad
S(P)=S.
\end{equation*}
The initial label~$\mathsf X(P)$ records whether the interaction at~$s(P)$ is
incoming or outgoing:
\begin{equation*}
\mathsf X(P)=
\pw{\mathsf I}{i(P)\in S(P)}{\mathsf O}{i(P)\notin S(P)}.
\end{equation*}
The terminal label~$\mathsf Y(P)$ records the next successful interaction on
the same site line, with~$\mathsf E$ meaning that there is no such interaction:
\begin{equation*}
\mathsf Y(P)=
\pwthree
{\mathsf I}{e(P)<\infty\text{ and }e(P)\in\mathcal T^{\mathsf I}\bb{i(P)}}
{\mathsf O}{e(P)<\infty\text{ and }e(P)\in\mathcal T^{\mathsf O}\bb{i(P)}}
{\mathsf E}{e(P)=\infty}.
\end{equation*}
\end{definition}

Let~$\mathcal P$ be the family of all patches. For this infinite-horizon family,
a patch with terminal label~$\mathsf E$ is infinite, while a patch with terminal
label~$\mathsf I$ or~$\mathsf O$ is finite.

For the representation at a finite time~$t$, patches which continue beyond~$t$
must be cut at the observation time. For~$P\in\mathcal P$ with~$s(P)\leq t$,
let~$P^{\downarrow t}$ be obtained by replacing its endpoint by
$\min\cb{e(P),t}$ and, when~$t<e(P)$, replacing its terminal label by
$\mathsf E$. Thus
\begin{equation*}
P^{\downarrow t}
=
\bb{
 i(P),
 [s(P),\min\cb{e(P),t}),
 \mathsf X(P)\mathsf Y^{\downarrow t}(P),
 S(P)
},
\end{equation*}
where
\begin{equation*}
\mathsf Y^{\downarrow t}(P)
=
\pw{\mathsf Y(P)}{e(P)\leq t}{\mathsf E}{t<e(P)}.
\end{equation*}
The families of bulk patches, end patches, and all finite-horizon patches are
\begin{align*}
\mathcal B_t
&=
\cb{P\in\mathcal P:e(P)\leq t},
\\
\mathcal E_t
&=
\cb{P^{\downarrow t}:P\in\mathcal P,\ s(P)\leq t<e(P)},
\\
\mathcal P_t
&=
\mathcal B_t\cup\mathcal E_t.
\end{align*}
Thus the bulk patches are completed by time~$t$, while the end patches are the
patches crossing~$t$, truncated at the observation time. The family~$\mathcal
P_t$ is finite almost surely. Distinct end patches have distinct sites, since
the successful-interaction times divide each site line into consecutive
intervals.

Figure~\ref{fig:patch-geometry} shows the full graphical realization and the
finite-horizon patch family induced by its successful-interaction skeleton.
\begin{figure}[htbp]
\centering
\resizebox{\linewidth}{!}{\begingroup
\definecolor{patchfigred}{RGB}{175,32,38}
\definecolor{patchfigactive}{gray}{0.86}
\definecolor{patchfigused}{gray}{0.32}
\definecolor{patchfigunused}{gray}{0.50}
\definecolor{patchfigoutline}{gray}{0.15}
\begin{tikzpicture}[x=0.92cm,y=0.78cm]
\tikzset{
  patchfig site/.style={draw=patchfigoutline,line width=0.48pt},
  patchfig active/.style={fill=patchfigactive,draw=none},
  patchfig success/.style={draw=black,line width=0.86pt,-{Stealth[length=2.0mm,width=1.35mm]}},
  patchfig unused/.style={draw=patchfigunused,line width=0.68pt,-{Stealth[length=1.75mm,width=1.18mm]}},
  patchfig horizon/.style={draw=patchfigoutline,line width=0.72pt},
  patchfig boundary/.style={draw=patchfigoutline,line width=0.40pt},
  patchfig label/.style={font=\scriptsize,text=black,inner sep=0pt},
  patchfig separator/.style={draw=patchfigred,line width=2.70pt,line cap=butt,line join=miter},
  patchfig small/.style={font=\scriptsize,text=patchfigoutline},
  patchfig used death/.style={font=\normalsize,text=patchfigused,inner sep=0pt},
  patchfig unused death/.style={font=\normalsize,text=patchfigunused,inner sep=0pt}
}
\newcommand{\PatchFigGrid}{%
  \foreach \x in {0,...,7}{\draw[patchfig site](\x,0)--(\x,8.4);}
  \draw[patchfig horizon](0,0)--(7,0);
  \draw[patchfig horizon](0,8.4)--(7,8.4);}
\newcommand{\PatchFigBarrier}[3]{%
  \begin{scope}
    \clip (#1,#2) rectangle ({#1+1},#3);
    \foreach \k in {-1,...,53}{%
      \pgfmathsetmacro{\yy}{-0.045+0.165*\k}
      \fill[patchfigred] (#1,\yy)--(#1,{\yy+0.075})--({#1+0.105},{\yy+0.155})--({#1+0.105},{\yy+0.080})--cycle;
      \fill[patchfigred] ({#1+1},\yy)--({#1+1},{\yy+0.075})--({#1+0.895},{\yy+0.155})--({#1+0.895},{\yy+0.080})--cycle;}
  \end{scope}}
\newcommand{\PatchFigInitial}{%
  \node[font=\small,inner sep=0pt](patchfiginf) at (3.5,-0.92){$\infty$};
  \draw[patchfig success]([xshift=-1.4mm]patchfiginf.north west)..controls(2.95,-0.40)and(2.05,-0.31)..(1.5,0);
  \draw[patchfig success](patchfiginf.north)--(3.5,0);
  \draw[patchfig success]([xshift=1.4mm]patchfiginf.north east)..controls(4.05,-0.40)and(4.95,-0.31)..(5.5,0);}
\begin{scope}
  \path[patchfig active] (0,1.1) rectangle (1,1.65);
  \path[patchfig active] (0,3.1) rectangle (1,8.4);
  \path[patchfig active] (1,0) rectangle (2,8.4);
  \path[patchfig active] (2,1.1) rectangle (3,5.7);
  \path[patchfig active] (2,7.85) rectangle (3,8.4);
  \path[patchfig active] (3,0) rectangle (4,2.1);
  \path[patchfig active] (3,4.1) rectangle (4,8.4);
  \path[patchfig active] (4,2.1) rectangle (5,6.1);
  \path[patchfig active] (4,7.1) rectangle (5,8.4);
  \path[patchfig active] (5,0) rectangle (6,8.4);
  \path[patchfig active] (6,7.1) rectangle (7,7.55);
  \PatchFigGrid
  \draw[patchfig success] (1.5,1.1)--(0.5,1.1);
  \draw[patchfig success] (1.5,1.1)--(2.5,1.1);
  \fill[black] (1.5,1.1) circle[radius=1.25pt];
  \draw[patchfig success] (3.5,2.1)--(2.5,2.1);
  \draw[patchfig success] (3.5,2.1)--(4.5,2.1);
  \fill[black] (3.5,2.1) circle[radius=1.25pt];
  \draw[patchfig success] (1.5,3.1)--(0.5,3.1);
  \draw[patchfig success] (1.5,3.1)--(2.5,3.1);
  \fill[black] (1.5,3.1) circle[radius=1.25pt];
  \draw[patchfig success] (2.5,4.1)--(1.5,4.1);
  \draw[patchfig success] (2.5,4.1)--(3.5,4.1);
  \fill[black] (2.5,4.1) circle[radius=1.25pt];
  \draw[patchfig success] (3.5,5.1)--(2.5,5.1);
  \draw[patchfig success] (3.5,5.1)--(4.5,5.1);
  \fill[black] (3.5,5.1) circle[radius=1.25pt];
  \draw[patchfig success] (4.5,6.1)--(3.5,6.1);
  \draw[patchfig success] (4.5,6.1)--(5.5,6.1);
  \fill[black] (4.5,6.1) circle[radius=1.25pt];
  \draw[patchfig success] (5.5,7.1)--(4.5,7.1);
  \draw[patchfig success] (5.5,7.1)--(6.5,7.1);
  \fill[black] (5.5,7.1) circle[radius=1.25pt];
  \draw[patchfig success] (1.5,7.85)--(0.5,7.85);
  \draw[patchfig success] (1.5,7.85)--(2.5,7.85);
  \fill[black] (1.5,7.85) circle[radius=1.25pt];
  \draw[patchfig unused] (4.5,0.55)--(3.5,0.55);
  \draw[patchfig unused] (4.5,0.55)--(5.5,0.55);
  \node[patchfig used death] at (0.5,1.65) {$\times$};
  \node[patchfig used death] at (2.5,5.7) {$\times$};
  \node[patchfig used death] at (6.5,7.55) {$\times$};
  \node[patchfig unused death] at (6.5,2.65) {$\times$};
  \PatchFigInitial
  \node[patchfig small,anchor=west] at (7.08,8.4){$t$};
  \node[patchfig small,anchor=west] at (7.08,0){$0$};
\end{scope}
\begin{scope}[shift={(8.10,0)}]
  \path[patchfig active] (1,0) rectangle (2,1.1);
  \path[patchfig active] (1,1.1) rectangle (2,3.1);
  \path[patchfig active] (1,4.1) rectangle (2,7.85);
  \path[patchfig active] (2,3.1) rectangle (3,4.1);
  \path[patchfig active] (3,0) rectangle (4,2.1);
  \path[patchfig active] (3,4.1) rectangle (4,5.1);
  \path[patchfig active] (4,5.1) rectangle (5,6.1);
  \path[patchfig active] (5,6.1) rectangle (6,7.1);
  \PatchFigBarrier{0}{1.1}{3.1}
  \PatchFigBarrier{0}{3.1}{7.85}
  \PatchFigBarrier{0}{7.85}{8.4}
  \PatchFigBarrier{1}{0}{1.1}
  \PatchFigBarrier{1}{1.1}{3.1}
  \PatchFigBarrier{1}{3.1}{4.1}
  \PatchFigBarrier{1}{4.1}{7.85}
  \PatchFigBarrier{1}{7.85}{8.4}
  \PatchFigBarrier{2}{1.1}{2.1}
  \PatchFigBarrier{2}{2.1}{3.1}
  \PatchFigBarrier{2}{3.1}{4.1}
  \PatchFigBarrier{2}{4.1}{5.1}
  \PatchFigBarrier{2}{5.1}{7.85}
  \PatchFigBarrier{2}{7.85}{8.4}
  \PatchFigBarrier{3}{0}{2.1}
  \PatchFigBarrier{3}{2.1}{4.1}
  \PatchFigBarrier{3}{4.1}{5.1}
  \PatchFigBarrier{3}{5.1}{6.1}
  \PatchFigBarrier{3}{6.1}{8.4}
  \PatchFigBarrier{4}{2.1}{5.1}
  \PatchFigBarrier{4}{5.1}{6.1}
  \PatchFigBarrier{4}{6.1}{7.1}
  \PatchFigBarrier{4}{7.1}{8.4}
  \PatchFigBarrier{5}{0}{6.1}
  \PatchFigBarrier{5}{6.1}{7.1}
  \PatchFigBarrier{5}{7.1}{8.4}
  \PatchFigBarrier{6}{7.1}{8.4}
  \PatchFigGrid
  \draw[patchfig boundary] (0,1.1)--(1,1.1);
  \draw[patchfig boundary] (0,3.1)--(1,3.1);
  \draw[patchfig boundary] (0,7.85)--(1,7.85);
  \draw[patchfig boundary] (0,8.4)--(1,8.4);
  \draw[patchfig boundary] (1,0)--(2,0);
  \draw[patchfig boundary] (1,1.1)--(2,1.1);
  \draw[patchfig boundary] (1,3.1)--(2,3.1);
  \draw[patchfig boundary] (1,4.1)--(2,4.1);
  \draw[patchfig boundary] (1,7.85)--(2,7.85);
  \draw[patchfig boundary] (1,8.4)--(2,8.4);
  \draw[patchfig boundary] (2,1.1)--(3,1.1);
  \draw[patchfig boundary] (2,2.1)--(3,2.1);
  \draw[patchfig boundary] (2,3.1)--(3,3.1);
  \draw[patchfig boundary] (2,4.1)--(3,4.1);
  \draw[patchfig boundary] (2,5.1)--(3,5.1);
  \draw[patchfig boundary] (2,7.85)--(3,7.85);
  \draw[patchfig boundary] (2,8.4)--(3,8.4);
  \draw[patchfig boundary] (3,0)--(4,0);
  \draw[patchfig boundary] (3,2.1)--(4,2.1);
  \draw[patchfig boundary] (3,4.1)--(4,4.1);
  \draw[patchfig boundary] (3,5.1)--(4,5.1);
  \draw[patchfig boundary] (3,6.1)--(4,6.1);
  \draw[patchfig boundary] (3,8.4)--(4,8.4);
  \draw[patchfig boundary] (4,2.1)--(5,2.1);
  \draw[patchfig boundary] (4,5.1)--(5,5.1);
  \draw[patchfig boundary] (4,6.1)--(5,6.1);
  \draw[patchfig boundary] (4,7.1)--(5,7.1);
  \draw[patchfig boundary] (4,8.4)--(5,8.4);
  \draw[patchfig boundary] (5,0)--(6,0);
  \draw[patchfig boundary] (5,6.1)--(6,6.1);
  \draw[patchfig boundary] (5,7.1)--(6,7.1);
  \draw[patchfig boundary] (5,8.4)--(6,8.4);
  \draw[patchfig boundary] (6,7.1)--(7,7.1);
  \draw[patchfig boundary] (6,8.4)--(7,8.4);
  \node[patchfig label] at (0.5,2.1) {$\mathsf{II}$};
  \node[patchfig label] at (0.5,5.475) {$\mathsf{II}$};
  \node[patchfig label] at (0.5,8.125) {$\mathsf{IE}$};
  \node[patchfig label] at (1.5,0.55) {$\mathsf{IO}$};
  \node[patchfig label] at (1.5,2.1) {$\mathsf{OO}$};
  \node[patchfig label] at (1.5,3.6) {$\mathsf{OI}$};
  \node[patchfig label] at (1.5,5.975) {$\mathsf{IO}$};
  \node[patchfig label] at (1.5,8.125) {$\mathsf{OE}$};
  \node[patchfig label] at (2.5,1.6) {$\mathsf{II}$};
  \node[patchfig label] at (2.5,2.6) {$\mathsf{II}$};
  \node[patchfig label] at (2.5,3.6) {$\mathsf{IO}$};
  \node[patchfig label] at (2.5,4.6) {$\mathsf{OI}$};
  \node[patchfig label] at (2.5,6.475) {$\mathsf{II}$};
  \node[patchfig label] at (2.5,8.125) {$\mathsf{IE}$};
  \node[patchfig label] at (3.5,1.05) {$\mathsf{IO}$};
  \node[patchfig label] at (3.5,3.1) {$\mathsf{OI}$};
  \node[patchfig label] at (3.5,4.6) {$\mathsf{IO}$};
  \node[patchfig label] at (3.5,5.6) {$\mathsf{OI}$};
  \node[patchfig label] at (3.5,7.25) {$\mathsf{IE}$};
  \node[patchfig label] at (4.5,3.6) {$\mathsf{II}$};
  \node[patchfig label] at (4.5,5.6) {$\mathsf{IO}$};
  \node[patchfig label] at (4.5,6.6) {$\mathsf{OI}$};
  \node[patchfig label] at (4.5,7.75) {$\mathsf{IE}$};
  \node[patchfig label] at (5.5,3.05) {$\mathsf{II}$};
  \node[patchfig label] at (5.5,6.6) {$\mathsf{IO}$};
  \node[patchfig label] at (5.5,7.75) {$\mathsf{OE}$};
  \node[patchfig label] at (6.5,7.75) {$\mathsf{IE}$};
  \draw[patchfig separator] (0,7.85) -- (1,7.85) -- (2,7.85) -- (3,7.85) -- (3,6.1) -- (4,6.1) -- (4,7.1) -- (5,7.1) -- (6,7.1) -- (7,7.1);
  \draw[patchfig success] (1.5,1.1)--(0.5,1.1);
  \draw[patchfig success] (1.5,1.1)--(2.5,1.1);
  \fill[black] (1.5,1.1) circle[radius=1.25pt];
  \draw[patchfig success] (3.5,2.1)--(2.5,2.1);
  \draw[patchfig success] (3.5,2.1)--(4.5,2.1);
  \fill[black] (3.5,2.1) circle[radius=1.25pt];
  \draw[patchfig success] (1.5,3.1)--(0.5,3.1);
  \draw[patchfig success] (1.5,3.1)--(2.5,3.1);
  \fill[black] (1.5,3.1) circle[radius=1.25pt];
  \draw[patchfig success] (2.5,4.1)--(1.5,4.1);
  \draw[patchfig success] (2.5,4.1)--(3.5,4.1);
  \fill[black] (2.5,4.1) circle[radius=1.25pt];
  \draw[patchfig success] (3.5,5.1)--(2.5,5.1);
  \draw[patchfig success] (3.5,5.1)--(4.5,5.1);
  \fill[black] (3.5,5.1) circle[radius=1.25pt];
  \draw[patchfig success] (4.5,6.1)--(3.5,6.1);
  \draw[patchfig success] (4.5,6.1)--(5.5,6.1);
  \fill[black] (4.5,6.1) circle[radius=1.25pt];
  \draw[patchfig success] (5.5,7.1)--(4.5,7.1);
  \draw[patchfig success] (5.5,7.1)--(6.5,7.1);
  \fill[black] (5.5,7.1) circle[radius=1.25pt];
  \draw[patchfig success] (1.5,7.85)--(0.5,7.85);
  \draw[patchfig success] (1.5,7.85)--(2.5,7.85);
  \fill[black] (1.5,7.85) circle[radius=1.25pt];
  \PatchFigInitial
  \node[patchfig small,anchor=west] at (7.08,8.4){$t$};
  \node[patchfig small,anchor=west] at (7.08,0){$0$};
\end{scope}
\end{tikzpicture}
\endgroup
}
\caption{A realization of the graphical construction (left) and the induced
finite-horizon patch family at time~$t$ (right). On the left, gray shading marks
times when a site is active, black arrows are successful interactions, gray
arrows are rings at inactive sources, and crosses are deaths. On the right,
gray shading highlights patches with outgoing terminal boundary, which remain
active throughout. Red striped sides delimit patches, the labels give their
initial and terminal boundary types, and the thick red line separates bulk
patches from end patches.}
\label{fig:patch-geometry}
\end{figure}
\FloatBarrier

If~$P\in\mathcal E_t$ and~$u\geq t$ is finite, define its extension through
time~$u$ by
\begin{equation}
\label{eq:geometric-patch-extension}
P^{\uparrow u}
=
\bb{i(P),[s(P),u),\mathsf X(P)\mathsf E,S(P)}.
\end{equation}
Thus~$P^{\uparrow t}=P$.
